\documentclass[../DoAn.tex]{subfiles}
\begin{document}

\begin{center}
    \Large{\textbf{TÓM TẮT NỘI DUNG ĐỒ ÁN}}\\
\end{center}
\vspace{1cm}
Việc quản lý tài chính cá nhân là một nhu cầu cấp thiết, đặc biệt đối với giới trẻ – nhóm đối tượng thường có nguồn thu nhập hạn chế và chưa hình thành thói quen chi tiêu hợp lý, dẫn đến rủi ro mất kiểm soát tài chính. Hiện tại, nhiều người vẫn dựa vào các công cụ thủ công như sổ tay hoặc bảng tính Excel để theo dõi chi tiêu. Tuy nhiên, các phương pháp này bộc lộ nhiều hạn chế nghiêm trọng: chúng thiếu tính tự động trong việc ghi nhận, khó khăn trong tổng hợp dữ liệu, và không cung cấp được các phân tích trực quan cần thiết để người dùng có thể đánh giá và điều chỉnh hành vi tài chính của bản thân. Xuất phát từ thực tế này, đồ án tốt nghiệp lựa chọn hướng tiếp cận phát triển một Ứng dụng Quản lý Chi tiêu Cá nhân dưới dạng ứng dụng web/mobile, nhằm thay thế hoàn toàn các phương pháp thủ công và cung cấp khả năng xử lý dữ liệu tự động, tức thời. Giải pháp được xây dựng sử dụng cơ sở dữ liệu tập trung và giao diện thân thiện, tập trung vào các chức năng cốt lõi bao gồm: ghi nhận thu nhập và chi tiêu hằng ngày; thiết lập ngân sách theo các danh mục cụ thể; và tạo ra các báo cáo tài chính bằng biểu đồ trực quan. Hệ thống cho phép người dùng dễ dàng theo dõi dòng tiền, so sánh chi tiêu thực tế với ngân sách đã đặt ra, từ đó nhanh chóng nhận diện và khắc phục các khoản chi tiêu không hợp lý. Kết quả của đồ án là một hệ thống hoàn chỉnh và ổn định, giúp người dùng chủ động quản lý tài chính cá nhân hiệu quả hơn, góp phần hình thành thói quen chi tiêu có kỷ luật. Đóng góp chính của đề tài là xây dựng thành công một mô hình ứng dụng đơn giản, dễ sử dụng nhưng có tính mở rộng cao, sẵn sàng cho việc tích hợp các tính năng nâng cao trong tương lai, điển hình là gợi ý chi tiêu thông minh dựa trên dữ liệu lịch sử người dùng.
\begin{flushright}
Sinh viên thực hiện\\
\begin{tabular}{@{}c@{}}
\textit{(Ký và ghi rõ họ tên)}
\end{tabular}
\end{flushright}

\end{document}